\emph{(1).} $B$ is a sum of commuting unit bivectors of square $-1$ in orthogonal planes, and $v \cdot B$ turns the component of $v$ in each plane by a quarter turn: $e_{1} \cdot (e_{1}e_{2}) = e_{2}$, $e_{2} \cdot (e_{1}e_{2}) = -e_{1}$. So $J$ is orthogonal with $J^{2} = -1$, $v \cdot J(v) = 0$, and $J^{a}$, $a = 1, 2, 3$, has no eigenvalue $1$. The check verifies these in $\Cl(6,0)$.

\emph{(2).} Write $\R^{2n+2} = P_{0} \oplus \dots \oplus P_{n}$ with $P_{r}$ the plane of $e_{2r+1}, e_{2r+2}$, and $S^{2r+1}$ the unit sphere of $P_{0} \oplus \dots \oplus P_{r}$. Let $e_{2r}$ be the set of $v \in S^{2r+1}$ whose $P_{r}$-component is a non-negative multiple of $e_{2r+1}$, and $e_{2r+1}$ the set of $v \in S^{2r+1}$ whose $P_{r}$-component lies in the open quarter plane between $e_{2r+1}$ and $J(e_{2r+1})$. Then $e_{2r}$ is a closed $2r$-disc with boundary $S^{2r-1}$ (it is the join of $S^{2r-1}$ with the point $e_{2r+1}$), $e_{2r+1}$ is an open $(2r+1)$-cell (the join of $S^{2r-1}$ with an open arc) whose closure adds $e_{2r}$, $g e_{2r}$ and $S^{2r-1}$, the translates $g^{a}e_{k}$ are disjoint and cover $S^{2r+1} \setminus S^{2r-1}$, and $g$ permutes them freely. With orientations chosen inductively and carried by $g$, the boundary of $e_{2r+1}$ is $g e_{2r} - e_{2r}$ (the two ends of the arc), and the boundary of $e_{2r}$ is the whole sphere $S^{2r-1}$, the sum of its four top cells $g^{a}e_{2r-1}$. On the group ring $\Z[G]$, multiplication by $g - 1$ has kernel the multiples of $N$ (an element fixed by $g$ has equal coefficients) and $N u = (\text{sum of the coefficients of } u)\,N$ has kernel the elements of coefficient sum zero, which is the image of $g - 1$. So the complex is exact in degrees $0 < k < 2n+1$, as it must be for a sphere. The check verifies $\partial^{2} = 0$ and the two ranks.

\emph{(3).} \emph{Chain map.} The formulas for $\Delta_{p,q}$ and for $\partial$ depend on $p$ and $q$ only through their parities, apart from the vanishing of $\partial$ in degree zero. So the identity $\partial\Delta(e_{k}) = \Delta(\partial e_{k})$, with $\partial(u \otimes v) = \partial u \otimes v + (-1)^{p}u \otimes \partial v$, has components indexed by $(p, q)$ with $p + q = k - 1$ that depend only on the parities of $p$ and $q$ and on whether $p$ or $q$ is zero; all these cases occur for $k \leq 5$. The check verifies the identity for $k \leq 9$ in exact integer arithmetic, for groups of order $2$, $3$ and $4$, and that $\Delta(e_{0}) = e_{0} \otimes e_{0}$. (The underlying mechanism is $\sum_{r<j} g^{r} \otimes g^{j}$ intertwining $g \otimes g - 1 \otimes 1$ with $N \otimes g - 1 \otimes N$.) By the comparison theorem of (G), since $C \otimes C$ has no homology in degrees $0 < k < 2n+1$, $\Delta$ is chain homotopic through degree $2n$ to a cellular approximation of the geometric diagonal, so it computes cup products into degrees at most $2n$. For products into the top degree $2n+1$, include $S^{2n+1}$ in $S^{2n+3}$ as the subcomplex of cells of dimension at most $2n+1$: the restriction $H^{*}(\Sig_{n+1}) \to H^{*}(\Sig_{n})$ is a ring map which on cochains sends $\varepsilon_{k}$ to $\varepsilon_{k}$ for $k \leq 2n+1$, and in $\Sig_{n+1}$ the degree $2n+1$ is below the top.

\emph{Cochains.} An invariant cochain with trivial coefficients is a multiple of $\varepsilon_{k}$ in each degree. $(\delta\varepsilon_{k-1})(e_{k}) = \varepsilon_{k-1}(\partial e_{k})$ is $1 - 1 = 0$ for $k$ odd and $4$ for $k$ even. So the integral cochain complex is $\Z$ in each degree with coboundary $0, 4, 0, 4, \dots$: $H^{2r} = \Z_{4}$ for $1 \leq r \leq n$, $H^{2r-1} = 0$ for $1 \leq r \leq n$, $H^{2n+1} = \Z$. $(\varepsilon_{p} \cup \varepsilon_{q})(e_{p+q}) = (\varepsilon_{p} \otimes \varepsilon_{q})(\Delta_{p,q})$ is the number of terms of $\Delta_{p,q}$: $1$, or $\binom{4}{2} = 6$ when both are odd. So $[\varepsilon_{2}]^{r} = [\varepsilon_{2r}]$, a generator, and $y^{n+1} = 0$ for want of cells: $H^{\mathrm{even}} = \Z[y]/(4y, y^{n+1})$. Modulo $2$ all coboundaries vanish, $x^{2} = 6[\varepsilon_{2}] = 0$, $xu^{r} = [\varepsilon_{2r+1}]$, $u^{r} = [\varepsilon_{2r}]$, which is the stated ring; modulo $4$ all coboundaries vanish too and $[\varepsilon_{1}]^{2} = 2[\varepsilon_{2}]$. The sign conventions of the cup product change none of these statements. For the group of order two the same formulas give $\binom{2}{2} = 1$.

\emph{(4).} The integral cochain $\varepsilon_{1}$ lifts the mod-$2$ cocycle representing $x$, and $\delta\varepsilon_{1} = 4\varepsilon_{2} = 2 \cdot 2\varepsilon_{2}$, so $\beta x = [2\varepsilon_{2}] = 2y$, of order two in $\Z_{4}$. For the sequence $\Z_{2} \to \Z_{4} \to \Z_{2}$: every mod-$2$ cocycle $\varepsilon_{k}$ lifts to the mod-$4$ cochain $\varepsilon_{k}$, whose coboundary is $0$ or $4\varepsilon_{k+1} \equiv 0$ modulo $4$; so the lift is a cocycle and the Bockstein is zero.

\emph{(5).} $\shL$ is the flat line with holonomy $-1$ around the loop that is the image of the path $e_{1}$ from the base point to its translate; a section over the base point extends over that loop only with a sign change, so the obstruction cocycle is $\varepsilon_{1}$ modulo $2$: $w_{1}(\shL) = x$. For $\eta$: a section is a map $\sigma$ from the sphere to the plane with $\sigma(J(v)) = J\sigma(v)$; take $\sigma_{1}(v)$ the $P_{0}$-component of $v$. It is nowhere zero on the $1$-skeleton $S^{1}$, and on the boundary of $e_{2}$, which is $S^{1}$ traversed once (the sum $Ne_{1}$), it is the identity map of the circle, of degree $1$. The obstruction cocycle is therefore $\pm\varepsilon_{2}$ and the Euler class $\pm y$. For $\eta^{2}$: $\sigma_{2}(v) = w e_{1} w$ with $w$ the $P_{0}$-component of $v$ satisfies $\sigma_{2}(J(v)) = -\sigma_{2}(v)$, is nowhere zero on $S^{1}$, and is the double-angle map there (verified in the check at rational points of the circle), of degree $2$: Euler class $\pm 2y$. Tangent bundle: at $v \in S^{2n+1}$, $\R^{2n+2} = T_{v}S^{2n+1} \oplus \R v$, and the linear map $J$ acts on this splitting equivariantly, fixing the radial section $v \mapsto v$. So $TS^{2n+1} \oplus \underline{\R}$ is the product bundle with fibre $\R^{2n+2} = P_{0} \oplus \dots \oplus P_{n}$ on which $g$ acts by the quarter turn in each plane, and dividing by $G$ gives $T\Sig_{n} \oplus \underline{\R} \cong (n+1)\eta$.

\emph{(6).} $\chi^{4}$ is the flat line with trivial holonomy, so $\chi^{4} = 1$. By (K1), $s = \chi - 1 \in F^{2}$ with image $c_{1}(\eta) = \pm y$ by (5), and by multiplicativity $s^{r} \in F^{2r}$ with image $\pm y^{r}$, a generator of $H^{2r} = \Z_{4}$ by (3); $s^{n+1} \in F^{2n+2} = 0$. By (3) the odd cohomology is zero below the top degree and $\Z$ there, so (K1) applies and the reduced group has a filtration with quotients $\Z_{4}$, $n$ times: its order is $4^{n}$. If $\sum a_{r}s^{r} = \sum b_{r}s^{r}$ with $0 \leq a_{r}, b_{r} < 4$ and $j$ the least index with $a_{j} \neq b_{j}$, the difference lies in $F^{2j}$ with image $\pm(a_{j} - b_{j})y^{j} \neq 0$, a contradiction; so these $4^{n}$ elements are distinct and are all of the reduced group. Hence the ring map $\Z[\chi]/(\chi^{4} - 1, (\chi - 1)^{n+1}) \to K(\Sig_{n})$ is onto. The augmentation ideal of the presented ring is spanned by $s, \dots, s^{n}$ with relations $s^{j}\bigl((1+s)^{4} - 1\bigr) = s^{j}(4s + 6s^{2} + 4s^{3} + s^{4})$, $0 \leq j < n$, a triangular system with diagonal entries $4$: order $4^{n}$. A surjection between finite groups of the same order is a bijection. The order of $1 - \chi^{2} = -(2s + s^{2})$ is the least $N$ with $N(2s + s^{2})$ in the relation lattice, computed exactly in the check for $n \leq 9$.
